% =====================================================================
% chapters/03-methodology.tex
% PT: Capítulo 3 -- Metodologia. Exemplo de: figura desenhada com TikZ
%     (sem ficheiro de imagem), equação numerada, números e unidades com
%     siunitx e referências a figura, equação e secção.
% EN: Chapter 3 -- Methodology. Shows: a figure drawn with TikZ (no image
%     file needed), a numbered equation, numbers and units with siunitx,
%     and references to a figure, an equation and a section.
% =====================================================================
\chapter{Methodology}
\label{chap:methodology}

% PT: Parágrafo de abertura do capítulo.
% EN: Chapter opening paragraph.
This chapter shows how to include a diagram and an equation. In a real
dissertation, the methodology describes what was done and how, in enough
detail for someone else to repeat the work: the overall approach, the data,
the methods and the experimental setup.

% ---------------------------------------------------------------------
% PT: Exemplo em inglês
% EN: English example
% ---------------------------------------------------------------------
\section{English example}
\label{sec:method:en}

\subsection{Overview}
\label{sec:method:overview}

% PT: Figura TikZ. [htbp] deixa o LaTeX escolher o melhor sítio:
%     h = aqui, t = topo, b = fundo, p = página própria.
%     Não use \gls dentro de \caption (use \acrshort), porque a legenda
%     também é escrita na lista de figuras.
% EN: TikZ figure. [htbp] lets LaTeX choose the best place:
%     h = here, t = top, b = bottom, p = separate page.
%     Do not use \gls inside \caption (use \acrshort), because the caption
%     is also written to the list of figures.
\begin{figure}[htbp]
  \centering
  \begin{tikzpicture}[
      node distance=8mm,
      stage/.style={draw, rounded corners, fill=blue!8, align=center,
                    minimum height=12mm, text width=24mm, font=\small\sffamily},
      arrow/.style={-{Stealth[length=2.5mm]}, thick}
    ]
    \node[stage] (data) {Data\\collection};
    \node[stage, right=of data] (prep) {Pre-\\processing};
    \node[stage, right=of prep] (model) {Model\\training};
    \node[stage, right=of model] (eval) {Evaluation};
    \draw[arrow] (data) -- (prep);
    \draw[arrow] (prep) -- (model);
    \draw[arrow] (model) -- (eval);
    \draw[arrow, dashed] (eval.south) -- ++(0,-8mm)
      -| node[pos=0.25, below, font=\footnotesize\sffamily] {adjust and repeat}
      (model.south);
  \end{tikzpicture}
  \caption{Example of a methodology diagram drawn with TikZ: the stages of a
    generic \acrshort{ml} pipeline.}
  \label{fig:pipeline}
\end{figure}

% PT: \cref{fig:...} escreve "Figure 3.1" com ligação.
% EN: \cref{fig:...} writes "Figure 3.1" with a link.
\Cref{fig:pipeline} summarises the methodology. Data are collected and
pre-processed, a model is trained, and the model is evaluated; the dashed
arrow shows that training is repeated with new settings until the
evaluation criteria are met. Diagrams drawn with TikZ use the same fonts as
the text and scale without loss of quality.

\subsection{Model}
\label{sec:method:model}

% PT: Equação numerada: o número aparece à direita, por capítulo (3.1).
%     Refira-a com \cref{eq:...} ou \eqref{eq:...}.
% EN: Numbered equation: the number appears on the right, per chapter
%     (3.1). Refer to it with \cref{eq:...} or \eqref{eq:...}.
A \gls{cnn} classifier is trained by minimising the cross-entropy loss over
the $N$ training examples and $C$ classes, given in \cref{eq:loss}:
\begin{equation}
  \mathcal{L}(\theta) = -\frac{1}{N} \sum_{i=1}^{N} \sum_{c=1}^{C}
  y_{i,c} \log \hat{y}_{i,c}(\theta),
  \label{eq:loss}
\end{equation}
where $\theta$ are the model parameters, $y_{i,c}$ is $1$ if example $i$
belongs to class $c$ and $0$ otherwise, and $\hat{y}_{i,c}(\theta)$ is the
probability predicted by the model.

\subsection{Experimental Setup}
\label{sec:method:setup}

% PT: siunitx: \num{...} para números e \qty{número}{unidade} para
%     quantidades (o antigo \SI{...}{...} continua a funcionar).
% EN: siunitx: \num{...} for numbers and \qty{number}{unit} for quantities
%     (the older \SI{...}{...} still works).
The model was trained on \num{50000} images for \num{30} epochs, with a
learning rate of \num{1e-3} and a batch size of \num{64}. The experiments ran
on a workstation with a \qty{3.2}{\giga\hertz} processor, \qty{32}{\giga\byte}
of memory and a graphics card with \qty{12}{\giga\byte} of memory. The setup
must allow the evaluation in \cref{sec:results:en} to be reproduced.

% ---------------------------------------------------------------------
% PT: Exemplo em português
% EN: Portuguese example
% ---------------------------------------------------------------------
\section{Exemplo em português}
\label{sec:method:pt}

\begin{otherlanguage}{portuguese}

% PT: Referências em português: nome à mão + \ref; \eqref põe parênteses.
% EN: References in Portuguese: name by hand + \ref; \eqref adds parentheses.
A Figura~\ref{fig:pipeline} resume as etapas da metodologia, descritas na
Secção~\ref{sec:method:overview}. O modelo é treinado minimizando a função
de perda da Equação~\eqref{eq:loss}. Na metodologia, descreva o que foi feito
com detalhe suficiente para que outra pessoa possa repetir o trabalho:
dados, métodos, ferramentas e parâmetros.

\end{otherlanguage}
